% =====================================================================
%  WYKAZ SKRÓTÓW I OZNACZEŃ — treść z wspolne/slownik.md (sekcje 2 i 3.1).
%  Słownik jest źródłem prawdy: zmiana tam → zmiana tutaj.
% =====================================================================
\chapter*{Wykaz skrótów i~oznaczeń}   % mwrep sam dodaje go do spisu treści

{\small\setstretch{1.0}
\noindent\textbf{Skróty}\par\smallskip
\noindent
\begin{tabularx}{\linewidth}{@{}l>{\raggedright\arraybackslash}X@{}}
APXS    & spektrometr rentgenowski wzbudzany cząstkami alfa (\angielski{Alpha Particle X-ray Spectrometer})\\
CNT     & nanorurki węglowe (\angielski{Carbon Nanotubes})\\
CNT-PVP & kompozyt nanorurek węglowych i~poliwinylopirolidonu\\
DRIE    & głębokie trawienie reaktywne jonowe (\angielski{Deep Reactive Ion Etching})\\
EDS     & spektroskopia rentgenowska z~dyspersją energii (\angielski{Energy Dispersive X-ray Spectroscopy})\\
F--N    & Fowler--Nordheim (równanie, wykres)\\
FWHM    & szerokość połówkowa (\angielski{Full Width at Half Maximum})\\
KOH     & wodorotlenek potasu\\
MCA     & analizator wielokanałowy (\angielski{Multichannel Analyzer})\\
MEMS    & mikrosystemy elektromechaniczne (\angielski{Micro-Electro-Mechanical Systems})\\
NdFeB   & stop neodym--żelazo--bor (magnes trwały)\\
PIXL    & przyrząd łazika Perseverance (\angielski{Planetary Instrument for X-ray Lithochemistry})\\
PVP     & poliwinylopirolidon (\angielski{Polyvinylpyrrolidone})\\
RTG     & rentgenowski (tylko w~nazwie „źródło RTG MEMS”)\\
SDD     & krzemowy detektor dryfowy (\angielski{Silicon Drift Detector})\\
SEM     & elektronowy mikroskop skaningowy (\angielski{Scanning Electron Microscope})\\
Si(Li)  & detektor krzemowy kompensowany litem\\
TMAH    & wodorotlenek tetrametyloamoniowy (\angielski{Tetramethylammonium Hydroxide})\\
UV      & nadfiolet (\angielski{Ultraviolet})\\
WDS     & spektroskopia z~dyspersją długości fali (\angielski{Wavelength Dispersive X-ray Spectroscopy})\\
XRF     & fluorescencja rentgenowska (\angielski{X-ray Fluorescence})\\
ZAF     & poprawki na liczbę atomową (Z), absorpcję (A) i~fluorescencję (F)\\
\end{tabularx}

\bigskip
\noindent\textbf{Oznaczenia}\par\smallskip
\noindent
\begin{tabularx}{\linewidth}{@{}l>{\raggedright\arraybackslash}X>{\raggedright\arraybackslash}p{22mm}@{}}
$\UB$            & napięcie ekstrakcyjne (emiter--bramka) & \si{\volt}\\
$\UT$            & napięcie przyspieszające (emiter--target) & \si{\volt}, \si{\kilo\volt}\\
$\Ie$            & prąd emisji & \si{\micro\ampere}\\
$\IB$            & prąd bramki & \si{\micro\ampere}\\
$\IT$            & prąd targetu & \si{\micro\ampere}\\
$\RT$, $\RB$     & rezystory ograniczające w~obwodzie targetu i~bramki & \si{\kilo\ohm}\\
$p$              & ciśnienie & \si{\milli\bar}, \si{\pascal}\\
$L$              & jasność świecenia ekranu (analiza obrazu) & \si{\ju}\\
$\varphi$        & praca wyjścia & \si{\electronvolt}\\
$F$              & natężenie pola elektrycznego na powierzchni emitera & \si{\volt\per\metre}\\
$\beta$          & współczynnik wzmocnienia pola & --\\
$\Ezero$         & energia elektronów wiązki, $\Ezero = e\UT$ & \si{\kilo\electronvolt}\\
$\Ec$            & energia krytyczna (krawędź absorpcji) & \si{\kilo\electronvolt}\\
$\vec B$         & indukcja magnetyczna & \si{\tesla}\\
\multicolumn{3}{@{}l@{}}{\TODO{pozostałe symbole ze słownika (p.~3.2) po ustaleniu rozdz.~2}}\\
\end{tabularx}
\par}
